\begin{table}[!htbp]
\raggedright
\caption{Prevalence of flow moments under each classification rule.}
\label{tab:sup-prevalence}
\footnotesize
\begingroup
\setlength{\tabcolsep}{4pt}
\renewcommand{\arraystretch}{1.08}
\noindent
\begin{tabular*}{\textwidth}{@{\extracolsep{\fill}}lccccc@{}}
\toprule
Rule & Condition met & Gated flow & Median person & Person range & At zero \\
\midrule
A. Raw, both items $\geq$ 4 & 39.6 & 36.8 & 35.0 & 1.0 to 89.4 & 0/68 \\
A. Raw, both items $\geq$ 5 & 22.0 & 19.3 & 16.4 & 0.0 to 66.7 & 4/68 \\
B. Grand-mean $z$ & 30.6 & 27.8 & 26.7 & 0.0 to 75.0 & 2/68 \\
C. Within-person $z$ & 28.9 & 24.6 & 24.5 & 1.0 to 53.2 & 0/68 \\
D. Person-mean centred & 28.9 & 24.6 & 24.5 & 1.0 to 53.2 & 0/68 \\
\bottomrule
\end{tabular*}
\par\vspace{3pt}\noindent\parbox{\textwidth}{\footnotesize\textit{Note.} All shares are percentages of moments, except the last column, which counts participants who never reach flow. The gated criterion requires both the challenge-skill condition and a flow experience above the same cut. Relative rules force a within-person distribution and therefore assign above-average moments to every participant, which is why only the absolute rules can represent the absence of flow and why every prevalence claim in the text is made on the absolute metric.}
\endgroup
\end{table}
